\documentclass[12pt,a4paper]{article}
\usepackage[spanish,es-nodecimaldot]{babel}
\usepackage[utf8]{inputenc}
\usepackage[T1]{fontenc}
\usepackage{lmodern}
\usepackage[top=1.1cm,bottom=1.1cm,left=1.2cm,right=1.2cm]{geometry}
\usepackage{amsmath,amssymb}
\usepackage[table]{xcolor}
\usepackage{booktabs,array}
\usepackage{enumitem}

\setlength{\parindent}{0pt}
\setlength{\parskip}{10pt}
\pagestyle{empty}

\newcommand{\cuerpo}{\fontsize{14}{17}\selectfont}

\AtBeginDocument{%
  \setlength{\abovedisplayskip}{8pt}%
  \setlength{\belowdisplayskip}{8pt}%
  \setlength{\abovedisplayshortskip}{4pt}%
  \setlength{\belowdisplayshortskip}{4pt}%
  \cuerpo
}

\begin{document}

{\fontsize{17}{21}\selectfont\bfseries Diseño en bloques aleatorios\par}

\hspace{14pt}\begin{minipage}{0.95\linewidth}
El experimentador tiene a disposicion mediciones relativas a ``a'' tratamientos distribuidos en b bloques. Vamos a considerar el caso en q hay una observacion en cada tratamiento de cada bloque.
\end{minipage}

\begin{center}
{\fontsize{11.5}{14}\selectfont
\setlength{\tabcolsep}{5pt}
\begin{tabular}{l | c c c c | c | c | c}
\toprule
\textbf{Tratamiento} & \textbf{Bloque 1} & \textbf{Bloque 2} & $\ldots$ & \textbf{Bloque $b$} & \textbf{Total ($T_{i.}$)} & \textbf{Media ($\bar{y}_{i.}$)} & \textbf{Efecto ($\alpha_i$)} \\
\midrule
Tratamiento 1 & $y_{11}$ & $y_{12}$ & $\ldots$ & $y_{1b}$ & $T_{1.}$ & $\bar{y}_{1.}$ & $\alpha_1=\bar{y}_{1.}-\bar{y}_{..}$ \\
Tratamiento 2 & $y_{21}$ & $y_{22}$ & $\ldots$ & $y_{2b}$ & $T_{2.}$ & $\bar{y}_{2.}$ & $\alpha_2=\bar{y}_{2.}-\bar{y}_{..}$ \\
$\ldots$ & $\ldots$ & $\ldots$ & $\ldots$ & $\ldots$ & $\ldots$ & $\ldots$ & $\ldots$ \\
Tratamiento $a$ & $y_{a1}$ & $y_{a2}$ & $\ldots$ & $y_{ab}$ & $T_{a.}$ & $\bar{y}_{a.}$ & $\alpha_a=\bar{y}_{a.}-\bar{y}_{..}$ \\
\midrule
\textbf{Total Bloque} & $T_{.1}$ & $T_{.2}$ & $\ldots$ & $T_{.b}$ & $T_{..}$ & & \\
\textbf{Media Bloque} & $\bar{y}_{.1}$ & $\bar{y}_{.2}$ & $\ldots$ & $\bar{y}_{.b}$ & & $\bar{y}_{..}$ & \\
\textbf{Efecto Bloque} & $\beta_1$ & $\beta_2$ & $\ldots$ & $\beta_b$ & & & \\
\bottomrule
\end{tabular}
}
\end{center}

\hspace{14pt}El modelo matemático aditivo sin interacción es:
\[
y_{ij} = \mu + \alpha_i + \beta_j + \epsilon_{ij} \qquad \text{para } i=1,\ldots,a; \quad j=1,\ldots,b
\]

\hspace{14pt}Con las restricciones $\sum_{i=1}^{a}\alpha_i=0$ y $\sum_{j=1}^{b}\beta_j=0$, donde:

\begin{itemize}[leftmargin=40pt,itemsep=6pt,topsep=3pt,parsep=0pt]
  \item $\mu$: Gran media general.
  \item $\alpha_i=\bar{y}_{i.}-\bar{y}_{..}$: Efecto fijo del $i$-ésimo tratamiento.
  \item $\beta_j=\bar{y}_{.j}-\bar{y}_{..}$: Efecto fijo del $j$-ésimo bloque.
  \item $\epsilon_{ij} \sim iid\ N(0,\sigma^2)$: Error aleatorio.
\end{itemize}

\begin{align*}
C &= \frac{T_{..}^2}{a\cdot b} \\
SST &= \sum_{i=1}^{a}\sum_{j=1}^{b} y_{ij}^2 - C \\
SS(Tr) &= \frac{1}{b}\sum_{i=1}^{a} T_{i.}^2 - C \qquad \text{con } a-1 \text{ gl} \\
SS(Bl) &= \frac{1}{a}\sum_{j=1}^{b} T_{.j}^2 - C \qquad \text{con } b-1 \text{ gl} \\
SSE &= SST - SS(Tr) - SS(Bl) \qquad \text{con } (a-1)(b-1) \text{ gl}
\end{align*}

\begin{center}
{\fontsize{12}{14.5}\selectfont
\setlength{\tabcolsep}{5pt}
\begin{tabular}{|l|c|c|l|l|}
\hline
\textbf{Fuentes de} & \textbf{Grados de} & \textbf{Suma de} & \textbf{Media Cuadrada} & \textbf{F} \\
\textbf{Variación} & \textbf{Libertad} & \textbf{Cuadrados} & & \\
\hline
\textbf{Tratamientos} & a-1 & SS(Tr) & MS(Tr)=SS(Tr)/(a-1) & $F_{Tr} = \frac{MS(Tr)}{MSE}$ \\
\hline
\textbf{Bloques} & b-1 & SS(Bl) & MS(Bl)=SS(Bl)/(b-1) & $F_{Bl} = \frac{MS(Bl)}{MSE}$ \\
\hline
\textbf{Error} & (a-1).(b-1) & SSE & MSE=SSE/(a-1).(b-1) & \\
\hline
\textbf{Total} & a.b-1 & SST & & \\
\hline
\end{tabular}
}
\end{center}

\hspace{40pt}\begin{minipage}{0.8\linewidth}
para llegar al error hicimos:\\
Error = total-tratamientos-bloques
\end{minipage}

\end{document}
